\section*{Bab \ref{ch:g8:solids} --- Limas dan Kerucut}

\begin{solution}{exo:g8:solids:1}
Bidang empat: $4$ sisi, $6$ \omterm{def:g5:solids:def}{rusuk}, $4$ \omterm{def:g5:solids:def}{titik sudut}
($4 + 4 = 6 + 2$). Alas bersegi enam: $7$ sisi, $12$ \omterm{def:g5:solids:def}{rusuk}, $7$
\omterm{def:g5:solids:def}{titik sudut} ($7 + 7 = 12 + 2$).
\end{solution}

\begin{solution}{exo:g8:solids:2}
$V = \frac13 \times 36 \times 10 = 120$ cm$^3$.

Alas persegi panjang: $B = 24$ cm$^2$, jadi
$V = \frac13 \times 24 \times 7.5 = 60$ cm$^3$.
\end{solution}

\begin{solution}{exo:g8:solids:3}
$V = \frac13 \pi \times 25 \times 9 = 75\pi \approx 236$ cm$^3$.
\end{solution}

\begin{solution}{exo:g8:solids:4}
\omterm{def:g7:areas:prism}{Tabung}: $\pi \times 16 \times 12 = 192\pi \approx 603$ cm$^3$.
\omterm{def:g8:solids:def}{Kerucut}: $\frac{192\pi}{3} = 64\pi \approx 201$ cm$^3$.
Perbandingannya: \omterm{def:g8:solids:def}{kerucutnya} sepertiga \omterm{def:g7:areas:prism}{tabungnya}.
\end{solution}

\begin{solution}{exo:g8:solids:5}
$56 = \frac13 \times B \times 8$, jadi $B = \frac{56 \times 3}{8} = 21$
cm$^2$.
\end{solution}

\begin{solution}{exo:g8:solids:6}
Selimutnya terbuka menjadi juring cakram yang \omterm{def:g3:shapes:shapes}{jari-jarinya} adalah
\emph{garis pelukisnya} $s$ (yaitu jarak dari puncaknya ke tepinya ---
\omterm{def:g6:lines:objects}{ruas garis} itu terletak pada permukaannya), bukan tingginya $h$.
\end{solution}

\begin{solution}{exo:g8:solids:7}
$V = \frac13 \times 35^2 \times 22 = \frac13 \times 1225 \times 22
= \frac{26\,950}{3} \approx 8\,983$ m$^3$ --- kira-kira $9\,000$ meter
kubik.
\end{solution}

\begin{solution}{exo:g8:solids:8}
Tingginya: $h = \sqrt{13^2 - 5^2} = \sqrt{169 - 25} = \sqrt{144} = 12$
cm. \omterm{def:g6:measure:volume}{Volumenya}: $V = \frac13 \pi \times 25 \times 12 = 100\pi$ cm$^3$.
\end{solution}

\begin{solution}{exo:g8:solids:9}
\omterm{def:g3:shapes:shapes}{Jari-jarinya} sama, jadi \omterm{def:g6:measure:volume}{volume} \omterm{def:g8:solids:def}{kerucutnya} sepertiga \omterm{def:g6:measure:volume}{volume} \omterm{def:g7:areas:prism}{tabungnya}
\emph{untuk tinggi yang sama}: cairannya mencapai
$\frac{9}{3} = 3$ cm di dalam \omterm{def:g7:areas:prism}{tabungnya}.
\end{solution}

\begin{solution}{exo:g8:solids:10}
Pasirnya $=$ satu \omterm{def:g8:solids:def}{kerucut}: $V = \frac13 \pi \times 4 \times 4.5 = 6\pi
\approx 18.8$ cm$^3$. Jam pasirnya memuat dua \omterm{def:g8:solids:def}{kerucut} seperti itu,
jadi pasirnya persis mengisi setengah \omterm{def:g6:measure:volume}{volume} dalam seluruhnya.
\end{solution}

\begin{solution}{exo:g8:solids:11}
\emph{1.} Ya: rumus $V = \frac13 Bh$ berlaku untuk letak puncak yang
mana pun, asalkan $h$ adalah jarak \omterm{def:g6:lines:perp}{tegak lurus} ke bidang alasnya.
$V = \frac13 \times 36 \times 8 = 96$ cm$^3$.

\emph{2.} Diagonal alasnya: $\sqrt{6^2 + 6^2} = 6\sqrt2$ cm. \omterm{def:g5:solids:def}{Rusuk}
tegak yang terpanjang adalah sisi miring \omterm{def:g6:shapes:triangles}{segitiga siku-siku} dengan
sisi siku-siku $6\sqrt2$ (pada bidang alasnya) dan $8$ (tegak):
\[
\sqrt{(6\sqrt2)^2 + 8^2} = \sqrt{72 + 64} = \sqrt{136} \approx 11.7
\text{ cm}.
\]
\end{solution}

\begin{solution}{pb:g8:solids:1}
\textbf{1.} $G$ termasuk ke sisi atas $EFGH$ dan ke kedua
sisi samping $BCGF$ dan $DCGH$. Ketiga sisi yang \emph{tidak}
memuat $G$ adalah alasnya $ABCD$, sisi depannya $ABFE$ dan
sisi kirinya $ADHE$ (ketiganya berbagi \omterm{def:g5:solids:def}{titik sudut} $A$, yaitu pojok
yang di seberang $G$). Ketiga limasnya adalah $ABCDG$, $ABFEG$ dan
$ADHEG$: masing-masing punya sebuah sisi persegi kubusnya sebagai
alas dan \omterm{def:g5:solids:def}{titik sudut} jauhnya $G$ sebagai puncak. Masing-masingnya
\omterm{def:g8:solids:def}{limas} beralas persegi: $5$ sisi, $8$ \omterm{def:g5:solids:def}{rusuk}, $5$ \omterm{def:g5:solids:def}{titik sudut}
(\cref{ex:g8:solids:count}).

\textbf{2.} Diagonal panjangnya $[AG]$ menghubungkan satu-satunya dua
\omterm{def:g5:solids:def}{titik sudut} yang tidak termasuk ke satu pun dari ketiga alasnya.
Perputaran sepertiga putaran terhadap garis $(AG)$ mengirim kubusnya
ke dirinya sendiri (ketiga \omterm{def:g5:solids:def}{rusuk} yang berangkat dari $A$ --- menuju
$B$, $D$, $E$ --- terkocok dalam satu daur, dan begitu pula ketiga
\omterm{def:g5:solids:def}{rusuk} yang tiba di $G$). Perputaran itu membawa alas $ABCD$ ke $ADHE$, lalu
$ADHE$ ke $ABFE$, lalu kembali: puncaknya tetap di $G$, alasnya
berdaur. Jadi tiap limasnya terbawa ke \omterm{def:g8:solids:def}{limas} berikutnya: ketiganya
salinan yang identik.

\textbf{3.} Tiga keping identik yang mengisi kubusnya berbagi
\omterm{def:g6:measure:volume}{volumenya} sama rata:
\[
V_{\text{limas}} = \frac{a^3}{3} .
\]

\textbf{4.} \omterm{def:g8:solids:def}{Limas} $ABCDG$ punya alas $ABCD$, yang luasnya $a^2$,
dan puncaknya $G$ duduk \omterm{def:g6:lines:perp}{tegak lurus} di atas pojok $C$, pada tinggi
$CG = a$ (yaitu sebuah \omterm{def:g5:solids:def}{rusuk} tegak kubusnya). Seperti pada
\cref{exo:g8:solids:11}, rumusnya berlaku dengan
tinggi \omterm{def:g6:lines:perp}{tegak lurus} $h = a$:
\[
V = \frac13 \times a^2 \times a = \frac{a^3}{3} ,
\]
yang sesuai sempurna dengan pertanyaan 3 --- penguraiannya dan
rumusnya saling membenarkan.

\textbf{5.} Untuk $a = 6$:
$V = \frac{6^3}{3} = \frac{216}{3} = 72$ cm$^3$ masing-masing.
Percobaan menuang airnya adalah penguraian ini yang dicairkan:
\omterm{def:g7:areas:prism}{prisma} di atas alas $ABCD$ dengan tinggi $a$ adalah kubusnya sendiri,
dan ia memuat persis tiga kali isi limasnya --- di sini ketiga
isiannya bahkan tersusun secara geometris menjadi kubusnya.

\textbf{6.} Menghubungkan pusatnya $O$ ke keempat pojok tiap
sisinya menghasilkan enam \omterm{def:g8:solids:def}{limas}, satu untuk tiap sisi: alas persegi,
puncak $O$. Pusatnya terletak sama terhadap keenam sisinya (setiap
perputaran atau simetri kubusnya menahan $O$ lalu mempermutasikan
sisinya), jadi keenam limasnya identik. Tinggi masing-masingnya
adalah jarak dari $O$ ke sebuah sisi: yaitu setengah sisinya, $\frac{a}{2}$.

\textbf{7.} Enam keping identik berbagi kubusnya:
$V = \frac{a^3}{6}$ masing-masing.

\textbf{8.} Dengan rumusnya: luas alasnya $a^2$, tingginya $\frac a2$:
\[
V = \frac13 \times a^2 \times \frac{a}{2} = \frac{a^3}{6} .
\]
Kedua jawabannya sesuai.

\textbf{9.} Kedua penguraiannya menyangkut \omterm{def:g8:solids:def}{limas} dengan perbandingan
yang sangat istimewa yang dipahat dari sebuah \omterm{def:g5:solids:def}{kubus}; \omterm{def:g8:solids:def}{limas} yang
ramping, \omterm{def:g8:solids:def}{limas} yang miring sebelah, atau \omterm{def:g8:solids:def}{kerucut} tidak dapat disusun
dengan cara ini (tidak ada tiga \omterm{def:g8:solids:def}{kerucut} yang mengisi
sebuah \omterm{def:g7:areas:prism}{tabung}). Itulah sebabnya \cref{thm:g8:solids:volume}
\emph{diterima tanpa bukti} pada tingkat ini: penguraiannya dan
percobaan menuang airnya membuat $\frac13$ itu sepenuhnya masuk akal,
dan bukti umumnya yang jujur --- yaitu pengintegralan --- diberikan di
buku Sekolah Menengah.

\textbf{10.} $V = \frac{6^3}{6} = 36$ cm$^3$; dan menurut rumusnya,
$V = \frac13 \times 36 \times 3 = 36$ cm$^3$. Keduanya sesuai.

\textbf{11.} Pada segitiga $SAB$, titik tengah sisi
$[SA]$ dan $[SB]$ dihubungkan \omterm{def:g6:lines:objects}{ruas garis} yang \omterm{def:g6:lines:perp}{sejajar} dengan $(AB)$
dan panjangnya $\frac{AB}{2} = \frac{c}{2}$
(\cref{thm:g8:midpoints:direct}). Hal yang sama berlaku pada tiap
sisi tegaknya, jadi keempat titik tengahnya membentuk persegi bersisi
$\frac c2$, dengan sisinya \omterm{def:g6:lines:perp}{sejajar} dengan alasnya. Pada garis tegak
yang melalui puncaknya, potongannya lewat di titik tengah tingginya
(dengan teorema titik tengah lagi, pada segitiga yang melalui $S$,
pusat alasnya dan sebuah \omterm{def:g5:solids:def}{titik sudut} alasnya): jadi bidang
potongnya duduk pada tinggi $\frac h2$.

\textbf{12.} Keping atasnya adalah \omterm{def:g8:solids:def}{limas} beraturan beralas persegi
dengan sisi alas $\frac c2$ dan tinggi $\frac h2$ (alasnya adalah
potongannya, puncaknya adalah $S$). \omterm{def:g6:measure:volume}{Volumenya} adalah
\[
\frac13 \times \left(\frac c2\right)^2 \times \frac h2
= \frac13 \times \frac{c^2}{4} \times \frac h2
= \frac{1}{8} \times \frac{c^2 h}{3}
= \frac{V}{8} :
\]
yaitu seperdelapan aslinya --- bukan setengahnya.

\textbf{13.} \omterm{def:g8:solids:def}{Limas} terpancungnya menyimpan sisanya:
$V - \frac V8 = \frac{7V}{8}$. Potongan setengah tinggi itu membagi
\omterm{def:g6:measure:volume}{volumenya} dengan perbandingan $1 : 7$ --- lempeng bawah yang tampak
sederhana memuat tujuh kali \omterm{def:g6:measure:volume}{volume} puncaknya yang lancip.

\textbf{14.} Di bawah setengah tingginya terletak $\frac78$
\omterm{def:g6:measure:volume}{volumenya}. Dengan
$V = \frac13 \times 35^2 \times 22 = \frac{26\,950}{3} \approx
8\,983$ m$^3$ (\cref{exo:g8:solids:7}):
\[
\frac78 \times \frac{26\,950}{3} \approx 7\,860 \approx
7\,900 \text{ m}^3
\]
sampai ratusan terdekat --- tahap ``separuh bawah'' membersihkan
hampir delapan kali \omterm{def:g6:measure:volume}{volume} kaca tahap yang di atasnya.

\textbf{15.} Melipatduakan setiap ukurannya mengalikan luas alasnya
dengan $2^2 = 4$ dan tingginya dengan $2$, jadi \omterm{def:g6:measure:volume}{volumenya} dengan
$4 \times 2 = 8 = 2^3$. Menskalakan dengan $k$ mengalikan luas dengan
$k^2$ dan satu panjang lagi dengan $k$: jadi \omterm{def:g6:measure:volume}{volume} tumbuh dengan
faktor $k^3$. Keping atas pada pertanyaan 12 adalah salinan terskala
seluruh limasnya dengan $k = \frac12$, jadi \omterm{def:g6:measure:volume}{volumenya}
$\left(\frac12\right)^3 = \frac18$ seluruhnya --- itulah penjelasan
satu barisnya.

\end{solution}
